\section{Appendix A. Mathematical Derivations}

\subsection{A.1 Energy Threshold Derivation}

\textbf{Claim.} If $r^* = \min\{r : E(r) \geq \tau\}$ where

\begin{equation}
E(r) = \frac{\sum_{i=1}^{r} \sigma_i^2}{\sum_{j=1}^{\min(m,n)} \sigma_j^2}
\end{equation}

then the relative Frobenius error satisfies:

\begin{equation}
\frac{\|W - W_{r^*}\|_F}{\|W\|_F} \leq \sqrt{1 - \tau}
\end{equation}

\textbf{Proof.} By Eckart-Young:

\begin{equation}
\|W - W_{r^*}\|_F^2 = \sum_{i > r^*} \sigma_i^2
\end{equation}

And:

\begin{equation}
\|W\|_F^2 = \sum_{i=1}^{\min(m,n)} \sigma_i^2
\end{equation}

So:

\begin{equation}
\frac{\|W - W_{r^*}\|_F^2}{\|W\|_F^2} = 1 - E(r^*) \leq 1 - \tau
\end{equation}

Taking square roots gives (A.2). $\square$

\subsection{A.2 Symmetric Split Equivalence}

\textbf{Claim.} For any $\Sigma_r \succ 0$, the following are equal:

\begin{equation}
U_r \Sigma_r V_r^\top = (U_r \Sigma_r^{1/2})(\Sigma_r^{1/2} V_r^\top)
\end{equation}

\textbf{Proof.} Trivial by associativity. $\square$

\textbf{Numerical benefit.} Splitting $\Sigma^{1/2}$ between the factors bounds 
the magnitude of each factor by $\sqrt{\sigma_1}$, reducing FP16 overflow/
underflow risk.

\subsection{A.3 Parameter Reduction}

For a projection $W \in \mathbb{R}^{d \times d}$:

\begin{equation}
\text{Baseline} = d^2
\end{equation}

\begin{equation}
\text{Factorized} = 2 d r^* + 2 d r_a
\end{equation}

Ratio:

\begin{equation}
\frac{\text{Baseline}}{\text{Factorized}} = \frac{d^2}{2 d (r^* + r_a)} = \frac{d}{2(r^* + r_a)}
\end{equation}

For $d = 768$, $r^* = 128$, $r_a = 16$:

\begin{equation}
\frac{768}{2 \cdot 144} = 2.67
\end{equation}

\subsection{A.4 Residual Adapter Recovery}

The effective weight is:

\begin{equation}
W_\text{eff} = U_r V_r + B A
\end{equation}

We want $W_\text{eff} \approx W$. The SVD part captures $W_{r^*}$; the 
adapter captures the residual $W - W_{r^*}$:

\begin{equation}
W - W_\text{eff} = (W - W_{r^*}) - B A
\end{equation}

The adapter rank $r_a$ should be large enough to represent this residual.

\subsection{A.5 GQA Composition}

For GQA with $H_{kv} < H$:

\begin{equation}
\text{K projection output} = H_{kv} \cdot d_h
\end{equation}

Factorized K:

\begin{equation}
\text{K params} = 2 d r^* + 2 r_a (H_{kv} d_h)
\end{equation}

Since $H_{kv} d_h < d$, this gives additional savings.

\subsection{A.6 T5 Relative Bias}

T5 adds a relative position bias to attention scores:

\begin{equation}
\text{scores}_{ij} \leftarrow \text{scores}_{ij} + b_{\text{bucket}(i, j)}
\end{equation}

The bias is parameterized by an embedding of shape $(num\_buckets, num\_heads)$ 
and is \textbf{not modified} by our factorization.

\subsection{A.7 Causality Preservation}

For causal attention, we mask scores at positions $j > i$:

\begin{equation}
\text{mask}_{ij} = \begin{cases} 0 & j \leq i \\ -\infty & j > i \end{cases}
\end{equation}

The factorization does not affect causality because it only modifies the 
projection weights.

